% !PRINTED 690
\weberpage{690}

Jedes Differential $dJ$, dessen Untereck $\mathfrak{B}$ oder ein Teiler
von $\mathfrak{B}$ ist, kann dann in der Form dargestellt werden:
\weq{7}{dJ = c_1dJ_1 + c_2dJ_2 + c_3dJ_3 + \cdots}
mit konstanten Koeffizienten.

Die einfachsten Differentiale sind die, deren Untereck das
Nulleck ist, und deren Oberecke $\mathfrak{W}$ also von der Ordnung
$(2p - 2)$ sind. Diese heißen \gesperrt{Differentiale erster Gattung}
und werden mit $dW$ bezeichnet. Sie sind dadurch charakterisiert,
daß für jede Variable $z$ die Polygone $\mathfrak{U}^2\mathfrak{W}$ und $\mathfrak{Z}$ äquivalent sein
müssen, und die $\mathfrak{W}$ bilden also, ihre Existenz vorausgesetzt, eine
Polygonklasse $W$, deren Dimension zu bestimmen ist.

Das Polygon $\mathfrak{W}$ heißt das Grundpolygon des Differentials $dW$
und wird ein \gesperrt{vollständiges Polygon erster Gattung genannt}. Ist
$\mathfrak{W} = \mathfrak{A}\mathfrak{B}$, so heißen auch $\mathfrak{A}$ und $\mathfrak{B}$ Polygone erster
Gattung und zwar \gesperrt{Ergänzungspolygone} voneinander.

Jedes Polygon, das nicht Teiler eines vollständigen Polygons
erster Gattung ist, heißt \gesperrt{von der zweiten Gattung}.

Ist $\mathfrak{A}$ ein Polygon erster Gattung, so ist auch jedes mit $\mathfrak{A}$
äquivalente Polygon $\mathfrak{A}'$ von der ersten Gattung, und wir nennen
die Klasse $A$ von $\mathfrak{A}$ eine \gesperrt{Klasse erster Gattung}. Denn ist
$\mathfrak{A}\mathfrak{B} = \mathfrak{W}$, so ist $AB = W$, wenn $B$ die Klasse von $\mathfrak{B}$ ist, und
es ist $\mathfrak{A}'\mathfrak{B} = \mathfrak{W}'$ mit $\mathfrak{W}$ äquivalent, also von der ersten Gattung.

Ist also $q$ die Anzahl der linear unabhängigen Polygone
erster Gattung, die durch ein Polygon der Klasse $A$ teilbar sind,
so ist nach § 192, (14)
\weq{8}{q = (A,\ W) = (O,\ B),}
also gleich der Dimension der Ergänzungsklasse.

Ist $A$ eine Klasse zweiter Gattung, so gibt es kein durch
$\mathfrak{A}$ teilbares Polygon erster Gattung, und es ist
\[ (A,\ W) = 0. \]

\begin{center}
§ 198. \textbf{\gesperrt{Die Polygonschar erster Gattung.}}
\end{center}

Um die Dimension $(O, W)$ der Polygonschar erster Gattung zu
bestimmen und damit zugleich ihre Existenz nachzuweisen, nehmen
wir eine beliebige Variable $z$ in $\Omega$ mit dem Verzweigungseck
$\mathfrak{Z}$ und dem Untereck $\mathfrak{U}$. Ist $dW$ ein Differential erster Gattung,
so ist
\weq{1}{w = \frac{dW}{dz} = \frac{\mathfrak{W}\mathfrak{U}^2}{\mathfrak{Z}}}
%CONT>>
